\documentclass{ximera}
\title{Properties of Logarithms}

\outcome{Apply the product, quotient, and power rules of logarithms.}
\outcome{Expand logarithmic expressions using logarithm properties.}
\outcome{Condense logarithmic expressions into a single logarithm.}
\outcome{Evaluate logarithmic expressions using known values and properties.}
\outcome{Apply the change of base formula.}

\begin{document}
\begin{abstract}
\end{abstract}
\maketitle

\textbf{Learning Objectives:}
\begin{itemize}
  \item Apply the product, quotient, and power rules of logarithms.
  \item Expand logarithmic expressions completely.
  \item Condense logarithmic expressions into a single logarithm.
  \item Evaluate expressions using given logarithm values and properties.
  \item Apply the change of base formula: $\log_b x = \dfrac{\ln x}{\ln b}$.
\end{itemize}
\medskip

\noindent\textbf{Key Rules} ($A,B>0$; $b>0$, $b\neq1$):
\begin{itemize}
  \item \textbf{Product:} $\log_b(AB) = \log_b A + \log_b B$
  \item \textbf{Quotient:} $\log_b\!\left(\dfrac{A}{B}\right) = \log_b A - \log_b B$
  \item \textbf{Power:} $\log_b(A^C) = C\cdot\log_b A$
  \item \textbf{Change of Base:} $\log_b x = \dfrac{\ln x}{\ln b} = \dfrac{\log x}{\log b}$
\end{itemize}

\bigskip
\section*{Quick Check}

\begin{problem}
Expand and simplify completely (assume all expressions are defined):
$$\ln\!\left(\frac{x^{3}\sqrt{\,1+x\,}}{e^{2}(x-2)^{4}}\right).$$
\begin{hint}
Apply the quotient rule, then the product rule, then the power rule. Recall $\ln e^2 = 2$.
\end{hint}
\begin{multipleChoice}
  \choice[correct]{$3\ln x + \tfrac{1}{2}\ln(1+x) - 2 - 4\ln(x-2)$}
  \choice{$3\ln x + \ln(1+x) - 2 - 4\ln(x-2)$}
  \choice{$3\ln x + \tfrac{1}{2}\ln(1+x) + 2 - 4\ln(x-2)$}
  \choice{$\ln(x^{3}+ \sqrt{1+x}) - 2 - 4\ln(x-2)$}
  \choice{$3\ln x + \tfrac{1}{2}\ln(1+x) - \ln(2) - 4\ln(x-2)$}
\end{multipleChoice}
\end{problem}

\begin{problem}
Condense into a single logarithm:
$$2\log_{5}(x-1) - \tfrac{1}{2}\log_{5}(x+3) + 3.$$
\begin{hint}
Apply the power rule first: $2\log_5(x-1)=\log_5(x-1)^2$ and $\tfrac{1}{2}\log_5(x+3)=\log_5\sqrt{x+3}$. Then write $3 = \log_5 5^3 = \log_5 125$ and combine.
\end{hint}
\begin{multipleChoice}
  \choice{$\log_{5}\!\left(\dfrac{(x-1)^{2}}{\sqrt{x+3}}\right)$}
  \choice[correct]{$\log_{5}\!\left(125\,\dfrac{(x-1)^{2}}{\sqrt{x+3}}\right)$}
  \choice{$\log_{5}\!\left( \dfrac{(x-1)^{2}}{125\,(x+3)^{1/2}}\right)$}
  \choice{$\log_{5}\!\left( \dfrac{(x-1)}{(x+3)}\right)^{2} + 3$}
  \choice{$3\log_{5}\!\left( \dfrac{(x-1)^{2}}{(x+3)^{1/2}}\right)$}
\end{multipleChoice}
\end{problem}

\begin{problem}
Rewrite $\log_{4}(7)$ using natural logarithms.
\begin{hint}
Use the change of base formula: $\log_b x = \dfrac{\ln x}{\ln b}$.
\end{hint}
\begin{multipleChoice}
  \choice{$\dfrac{\ln 4}{\ln 7}$}
  \choice[correct]{$\dfrac{\ln 7}{\ln 4}$}
  \choice{$\ln(4\cdot 7)$}
  \choice{$\ln 7 - \ln 4$}
  \choice{$\ln 4 + \ln 7$}
\end{multipleChoice}
\end{problem}

\begin{problem}
Use properties of logarithms to expand:
$$\log_{10}\!\big( 5x^{3}(x+1)^{2} \big).$$
\begin{hint}
Use the product rule to split the three factors, then apply the power rule to each.
\end{hint}
\begin{multipleChoice}
  \choice[correct]{$\log 5 + 3\log x + 2\log(x+1)$}
  \choice{$\log(5x^{3}) + \log(x+1)^{2}$}
  \choice{$\log(5+x^{3}+(x+1)^{2})$}
  \choice{$3\log(5x(x+1))$}
  \choice{$\log 5 + \log x + \log(x+1)$}
\end{multipleChoice}
\end{problem}

\begin{problem}
Without a calculator, determine the relationship between $\ln(15)$ and $2\ln(6) - \ln(10)$.
\begin{hint}
Condense: $2\ln(6)-\ln(10) = \ln(36)-\ln(10) = \ln\!\left(\tfrac{36}{10}\right) = \ln(3.6)$. Compare $\ln(15)$ with $\ln(3.6)$ using the fact that $\ln$ is increasing.
\end{hint}
\begin{multipleChoice}
  \choice[correct]{$\ln(15) > 2\ln(6) - \ln(10)$}
  \choice{$\ln(15) = 2\ln(6) - \ln(10)$}
  \choice{$\ln(15) < 2\ln(6) - \ln(10)$}
  \choice{Cannot be determined without a calculator}
  \choice{It depends on the logarithm base}
\end{multipleChoice}
\end{problem}


\bigskip
\section*{Extended Practice}

\begin{problem}
Expand the following logarithmic expressions and simplify completely.
\begin{enumerate}
  \item[(a)] $\log_4 (7xy)$
  \medskip
  \item[(b)] $\log_7 \!\left(\dfrac{1}{7}mn^2\right)$
  \medskip
  \item[(c)] $\log_2\!\left(\dfrac{x^{10}}{yz}\right)$
  \medskip
  \item[(d)] $\log\!\left(\dfrac{\sqrt{d^2+1}}{10000}\right)$
  \medskip
  \item[(e)] $\ln\!\left(\dfrac{\sqrt[4]{pq}}{t^3m}\right)$
  \medskip
  \item[(f)] $\log_5\sqrt[3]{x\sqrt{5}}$
  \medskip
  \item[(g)] $\log\!\left[\dfrac{2x(x^2-4)^8}{\sqrt{4-3x}}\right]$
  \medskip
  \item[(h)] $\log_2\!\left[\dfrac{4a^2\sqrt{3-b}}{c(b+4)^2}\right]$
\end{enumerate}

\begin{solution}
\begin{enumerate}
  \item[(a)] $\log_4(7xy) = \log_4 7 + \log_4 x + \log_4 y$

  \item[(b)]
  \begin{align*}
    \log_7\!\left(\tfrac{1}{7}mn^2\right) &= \log_7 7^{-1} + \log_7 m + \log_7 n^2\\
                                           &= -1 + \log_7 m + 2\log_7 n
  \end{align*}

  \item[(c)]
  \begin{align*}
    \log_2\frac{x^{10}}{yz} &= 10\log_2 x - \log_2 y - \log_2 z
  \end{align*}

  \item[(d)]
  \begin{align*}
    \log\frac{\sqrt{d^2+1}}{10000} &= \tfrac{1}{2}\log(d^2+1) - 4
  \end{align*}

  \item[(e)]
  \begin{align*}
    \ln\frac{\sqrt[4]{pq}}{t^3m} &= \tfrac{1}{4}\ln p + \tfrac{1}{4}\ln q - 3\ln t - \ln m
  \end{align*}

  \item[(f)]
  \begin{align*}
    \log_5\sqrt[3]{x\sqrt{5}} &= \log_5\!\left(x\cdot 5^{1/2}\right)^{1/3} = \tfrac{1}{3}\log_5 x + \tfrac{1}{6}
  \end{align*}

  \item[(g)]
  \begin{align*}
    \log\frac{2x(x^2-4)^8}{\sqrt{4-3x}} &= \log 2 + \log x + 8\log(x-2) + 8\log(x+2) - \tfrac{1}{2}\log(4-3x)
  \end{align*}

  \item[(h)]
  \begin{align*}
    \log_2\frac{4a^2\sqrt{3-b}}{c(b+4)^2} &= 2 + 2\log_2 a + \tfrac{1}{2}\log_2(3-b) - \log_2 c - 2\log_2(b+4)
  \end{align*}
\end{enumerate}
\end{solution}
\end{problem}


\begin{problem}
Combine the following logarithmic expressions and simplify completely.
\begin{enumerate}
  \item[(a)] $\log 5 + \log p$
  \medskip
  \item[(b)] $\log_{15} 3 + \log_{15} 5$
  \medskip
  \item[(c)] $\log_7 98 - \log_7 2$
  \medskip
  \item[(d)] $\log 150 - \log 3 - \log 5$
  \medskip
  \item[(e)] $8\log_3 x - 2\log_3 z - 7\log_3 y$
  \medskip
  \item[(f)] $15\log c - \dfrac{1}{4}\log d - \dfrac{3}{4}\log k$
  \medskip
  \item[(g)] $\dfrac{1}{3}\!\left[ 12\ln(x-5)+\ln x - \ln x^3 \right]$
\end{enumerate}

\begin{solution}
\begin{enumerate}
  \item[(a)] $\log 5 + \log p = \log 5p$

  \item[(b)] $\log_{15} 3 + \log_{15} 5 = \log_{15} 15 = 1$

  \item[(c)] $\log_7 98 - \log_7 2 = \log_7 49 = 2$

  \item[(d)] $\log 150 - \log 3 - \log 5 = \log\dfrac{150}{15} = \log 10 = 1$

  \item[(e)] $8\log_3 x - 2\log_3 z - 7\log_3 y = \log_3\dfrac{x^8}{y^7z^2}$

  \item[(f)] $15\log c - \dfrac{1}{4}\log d - \dfrac{3}{4}\log k = \log\dfrac{c^{15}}{\sqrt[4]{dk^3}}$

  \item[(g)]
  \begin{align*}
    \tfrac{1}{3}\!\left[12\ln(x-5)+\ln x - \ln x^3\right] &= 4\ln(x-5) + \tfrac{1}{3}\ln x - \ln x\\
    &= \ln\frac{(x-5)^4\sqrt[3]{x}}{x} = \ln\frac{(x-5)^4}{\sqrt[3]{x^2}}
  \end{align*}
\end{enumerate}
\end{solution}
\end{problem}


\begin{problem}
Suppose $\log_3 x = 3$, $\log_3 y = 7$, and $\log_3 z = 10$. Evaluate $\ds\log_3\!\left(\dfrac{x^2y}{\sqrt{z}}\right)$.

\begin{solution}
\begin{align*}
  \log_3\frac{x^2y}{\sqrt{z}} &= 2(\log_3 x) + (\log_3 y) - \tfrac{1}{2}(\log_3 z)\\
  &= 2(3) + 7 - \tfrac{1}{2}(10) = 6 + 7 - 5 = \mathbf{8}
\end{align*}
\end{solution}
\end{problem}


\begin{problem}
Suppose $\log_b 2 = 0.356$, $\log_b 3 = 0.565$, and $\log_b 5 = 0.827$. Evaluate each expression.
\begin{enumerate}
  \item[(a)] $\log_b 10$
  \medskip
  \item[(b)] $\log_b 81$
  \medskip
  \item[(c)] $\log_b 225$
  \medskip
  \item[(d)] $\log_b \dfrac{6}{5}$
\end{enumerate}

\begin{solution}
\begin{enumerate}
  \item[(a)] $\log_b 10 = \log_b(2\cdot 5) = 0.356 + 0.827 = \mathbf{1.183}$

  \item[(b)] $\log_b 81 = \log_b 3^4 = 4(0.565) = \mathbf{2.26}$

  \item[(c)] $\log_b 225 = \log_b 15^2 = 2(\log_b 3 + \log_b 5) = 2(0.565+0.827) = \mathbf{2.784}$

  \item[(d)] $\log_b\dfrac{6}{5} = \log_b 2 + \log_b 3 - \log_b 5 = 0.356 + 0.565 - 0.827 = \mathbf{0.094}$
\end{enumerate}
\end{solution}
\end{problem}


\begin{problem}
Use the change of base formula to write $(\log_4 11)(\log_{11} 3)$ as a single logarithm.

\begin{solution}
Apply $\log_b x = \dfrac{\ln x}{\ln b}$:
$$(\log_4 11)(\log_{11} 3) = \frac{\ln 11}{\ln 4}\cdot\frac{\ln 3}{\ln 11} = \frac{\ln 3}{\ln 4} = \log_4 3$$
\end{solution}
\end{problem}

\end{document}
